% 4.1 章 · 抛物型方程与时间推进
% 这一章把「热方程怎么往前推」的全部结论挂在同一个对象上：放大因子 G(φ)。
%   θ 族格式 → G(φ) 的闭式 → 稳定域 → 有限网格阈值 → 时间误差首项 → 耗散方向 → 边界守恒。
% 每个定理后面都给本章断言里的**实测数字**，纸面结论与机器结论同源。
% 断言附表由 tools/build_book.py 从实际运行结果生成，\input 进来，不手抄。

\numlabchapter{抛物型方程与时间推进}{4.1}

\section{这一章要解决什么}

一维热传导方程
\begin{equation}
  \theta_t=\alpha\,\theta_{xx},\qquad x\in(0,L),\ t>0
  \label{eq:heat}
\end{equation}
是抛物型方程的原型：它只有耗散、没有色散，解只会越来越光滑。数值上要回答四个问题，
而它们的答案全部挂在同一个对象上——\textbf{放大因子} $G(\varphi)$。

\begin{enumerate}
  \item \textbf{怎么推？} 空间用二阶中心差分，时间用 $\theta$ 族：
        $\theta=0$ 是显式（FTCS）、$\theta=1$ 是全隐式（BTCS）、$\theta=\frac12$ 是
        Crank--Nicolson。三种格式共用同一个空间离散，差别只在时间。
  \item \textbf{稳定吗？} 把初值展开成 Fourier 模态，一个模态一个乘子：
        $|G(\varphi)|\le1$ 就是全部判据。显式给出 $r\le\frac12$，
        隐式与 CN 无条件稳定。但"$\frac12$"只在连续极限下才是一句完整的话
        （定理~\ref{thm:finite}）。
  \item \textbf{准吗？} 稳定与精确是两件事。要测时间阶，必须把空间误差**解析地扣掉**：
        单模态下数值解逐点等于 $G(\varphi_k)^n\sin(k\pi x/L)$（定理~\ref{thm:mode}），
        于是时间误差有闭式（推论~\ref{cor:timeerr}）。拿连续解当靶子测时间阶，
        CN 会被量成 $0$ 阶（实测 $-0.0004$）。
  \item \textbf{耗散在哪？} 短波掉得最快，这就是"数值解总比真解光滑"的原因；
        而耗散误差的**符号**由 $r$ 与 $\theta$ 共同决定，在 $\theta=0$ 时会随 $r$ 反转
        （定理~\ref{thm:symbol}）。
\end{enumerate}

记号：$r=\alpha\Delta t/\Delta x^2$ 是格点比，$\varphi=k\Delta x\in[0,\pi]$ 是无量纲波数
（$\pi$ 是网格能表示的最短波，Nyquist），$s^2=\sin^2(\varphi/2)$，
$u=\frac12\varepsilon_{\text{mach}}=2^{-53}$，$n$ 是格点数。
所有"实测"数字都来自 \texttt{tests/4.1/} 的断言报告，
最坏的一帧（$n_x=101$、$r=0.05$、$t=0.05$）构建耗时 $0.12$ 秒。

\section{$\theta$ 族格式与截断误差}

\begin{definition}[$\theta$ 族]
设 $u_j^n$ 是 $\theta(x_j,t^n)$ 的近似，$\tilde D$ 是以 $\Delta x^2$ 为单位的二阶中心差分
\begin{equation}
  (\tilde Du)_j=u_{j+1}-2u_j+u_{j-1},
\end{equation}
则 $\theta$ 族格式是
\begin{equation}
  \bigl(I-\theta r\tilde D\bigr)u^{n+1}=\bigl(I+(1-\theta)r\tilde D\bigr)u^n .
  \label{eq:theta}
\end{equation}
$\theta=0$ 时右端的 $I+r\tilde D$ 显式可用；$\theta>0$ 时每一步要解一个三对角系统。
\end{definition}

\begin{theorem}[截断误差的阶]
\label{thm:trunc}
设 $\theta$ 是~\eqref{eq:heat}的光滑解。把 $\theta$ 代入~\eqref{eq:theta}，局部截断误差为
\begin{equation}
  \tau=\Bigl(\frac12-\theta\Bigr)\Delta t\,\theta_{xxxx}+O(\Delta t^2)+O(\Delta x^2).
\end{equation}
因此 $\theta\ne\frac12$ 时格式是一阶，$\theta=\frac12$（CN）时是二阶。
\end{theorem}

\begin{proof}
记 $D=\partial_x^2$。先展开时间方向：
\begin{equation}
  \frac{\theta(t+\Delta t)-\theta(t)}{\Delta t}
  =\theta_t+\frac{\Delta t}{2}\theta_{tt}+\frac{\Delta t^2}{6}\theta_{ttt}+O(\Delta t^3).
\end{equation}
再展开空间方向。由 $\theta_t=\alpha\theta_{xx}$ 得
$\theta_{tt}=\alpha^2\theta_{xxxx}$，且 $\partial_t D\theta=D\theta_t=\alpha\theta_{xxxx}$，
于是
\begin{equation}
  D\theta(t+\Delta t)=D\theta(t)+\Delta t\,\alpha\theta_{xxxx}+O(\Delta t^2),
\end{equation}
从而
\begin{equation}
  \theta D\theta(t+\Delta t)+(1-\theta)D\theta(t)
  =\alpha\theta_{xx}+\theta\Delta t\,\alpha\theta_{xxxx}+O(\Delta t^2).
\end{equation}
两式相减（并用 $\theta_t=\alpha\theta_{xx}$）：
\begin{equation}
  \tau=\frac{\Delta t}{2}\alpha^2\theta_{xxxx}-\theta\Delta t\,\alpha\theta_{xxxx}+O(\Delta t^2)
  =\Delta t\,\alpha\theta_{xxxx}\Bigl(\frac{\alpha}{2}\cdot\frac{1}{\alpha}\Bigr)\Bigl(\frac12-\theta\Bigr)+O(\Delta t^2).
\end{equation}
取 $\alpha=1$ 即得结论。空间项的 $O(\Delta x^2)$ 来自 $\tilde D\theta=\Delta x^2(D\theta+O(\Delta x^2))$。
\end{proof}

\textbf{实测}：单模态时间误差的首项系数（定理~\ref{thm:mode}的推论）在一阶格式上是
$O(\Delta t)$、在 CN 上是 $O(\Delta t^2)$；$nx=41$、$r=0.4$、$t=0.05$ 时 CN 的时间误差
是 $1.53\times10^{-7}$，而 BTCS 是 $3.71\times10^{-4}$——同一 $\Delta t$ 下差 $2400$ 倍。

\section{von Neumann 分析：放大因子与稳定域}

\begin{theorem}[放大因子]
\label{thm:g}
把 $u_j^n=G^n e^{\mathrm i\varphi j}$ 代入~\eqref{eq:theta}，得到
\begin{equation}
  G_\theta(\varphi)=\frac{1-4(1-\theta)r\,s^2}{1+4\theta r\,s^2},
  \qquad s^2=\sin^2\frac\varphi2\in[0,1].
  \label{eq:G}
\end{equation}
\end{theorem}

\begin{proof}
$\tilde D e^{\mathrm i\varphi j}=(e^{\mathrm i\varphi}-2+e^{-\mathrm i\varphi})e^{\mathrm i\varphi j}
=-4\sin^2\frac\varphi2\,e^{\mathrm i\varphi j}$。代入~\eqref{eq:theta}并约去 $e^{\mathrm i\varphi j}$：
$(1+4\theta rs^2)G=1-4(1-\theta)rs^2$。
\end{proof}

\begin{theorem}[稳定判据]
\label{thm:stab}
$\max_{\varphi\in[0,\pi]}|G_\theta(\varphi)|\le1$ 当且仅当
\begin{equation}
  r\,(1-2\theta)\le\frac12 ,
\end{equation}
即 $\theta\ge\frac12$ 时无条件稳定，$\theta<\frac12$ 时要求
$r\le\frac{1}{2(1-2\theta)}$。特别地 FTCS（$\theta=0$）要求 $r\le\frac12$。
\end{theorem}

\begin{proof}
记 $a=(1-\theta)r\ge0$、$b=\theta r\ge0$，则 $G=\frac{1-4as^2}{1+4bs^2}$。
分母 $\ge1>0$，分子随 $s^2$ 单调减、分母随 $s^2$ 单调增，故 $|G|$ 的极值只在
$s^2=0$ 或 $s^2=1$ 取到。

$s^2=0$ 时 $G=1$。$s^2=1$ 时 $|G|\le1$ 等价于两个不等式：
\begin{equation}
  \underbrace{1-4a\le1}_{\text{自动成立，因 }a\ge0},
  \qquad
  \underbrace{1-4a\ge-(1+4b)}_{\text{即 }4(a-b)\le2}.
\end{equation}
第二式即 $2r(1-2\theta)\le1$。$\theta\ge\frac12$ 时左端 $\le0$ 自动成立；否则给出
$r\le\frac{1}{2(1-2\theta)}$。FTCS 取 $\theta=0$ 得 $r\le\frac12$。
\end{proof}

\begin{corollary}[稳定性只由两个数决定]
只需检查 $\varphi=0$（$G=1$）与 $\varphi=\pi$（$G=\frac{1-4(1-\theta)r}{1+4\theta r}$）。
\end{corollary}

\begin{corollary}[CN 的"稳定"退化]
\label{cor:cn}
CN 的 $G(\pi)=\frac{1-2r}{1+2r}\to-1\ (r\to\infty)$：它永不增长，但每一步翻一次符号。
实测谱半径：$r=50$ 时 $0.9801$、$r=500$ 时 $0.9980$——"无条件稳定"在大步长下
退化成"不衰减、只抖动"。
\end{corollary}

\textbf{与收敛的关系}：$\theta$ 族与~\eqref{eq:heat}相容（定理~\ref{thm:trunc}），
故由 Lax 等价定理，\textbf{相容 $+$ 稳定 $\iff$ 收敛}。
"$r\le\frac12$"因此在 $L^2$ 意义下是收敛的充要条件，不是经验法则。

\section{有限网格上的阈值：为什么不是恰好 $\tfrac12$}

\begin{theorem}[有限网格阈值]
\label{thm:finite}
设网格上可分辨的最大波数是 $\varphi_{\max}$，则 FTCS 的稳定阈值是
\begin{equation}
  r_*=\frac{1}{2\sin^2(\varphi_{\max}/2)} ,
  \label{eq:rcrit}
\end{equation}
其中 $\varphi_{\max}$ 由边界条件（以及网格奇偶性）决定：
\begin{equation}
  \varphi_{\max}=
  \begin{cases}
    \dfrac{(n-2)\pi}{n-1}<\pi, & \text{Dirichlet},\\[2mm]
    \pi, & \text{Neumann},\\[2mm]
    \dfrac{2\pi\lfloor n/2\rfloor}{n}\le\pi, & \text{周期（仅偶数 }n\text{ 取到 }\pi).
  \end{cases}
\end{equation}
\end{theorem}

\begin{proof}
由定理~\ref{thm:stab}的证明，$|G|$ 随 $s^2$ 单调减（对 FTCS，$G=1-4rs^2$，
在 $s^2=\frac{1}{4r}$ 处过零后继续下降），故最危险的是 $s^2$ 最大的那个模态。
$s^2\le1$ 只能在允许的模态集合里取：Dirichlet 的模态是
$\varphi_k=\frac{k\pi}{n-1}$（$k=1,\dots,n-2$），最大只到 $\frac{(n-2)\pi}{n-1}$；
Neumann 的模态是 $\varphi_k=\frac{k\pi}{n-1}$（$k=0,\dots,n-1$），含 $\pi$；
周期网格的模态是 $\frac{2\pi m}{n}$，而 $m$ 与 $n-m$ 是同一个物理模态，
故只需 $m\le\lfloor n/2\rfloor$，奇数 $n$ 时取不到 $\pi$。
$r_*$ 由 $1-4r_*s_{\max}^2=-1$ 解出。
\end{proof}

\begin{corollary}[Dirichlet 阈值的二阶偏差]
\label{cor:dev}
Dirichlet 网格上
\begin{equation}
  r_*=\frac{1}{2\cos^2\frac{\pi}{2(n-1)}}
  =\frac12+\frac{\pi^2}{8(n-1)^2}+O(n^{-4}).
\end{equation}
\end{corollary}

\begin{proof}
$\varphi_{\max}=\pi-\frac{\pi}{n-1}$，故 $\sin\frac{\varphi_{\max}}2=\cos\frac{\pi}{2(n-1)}$
$=1-\frac{\pi^2}{8(n-1)^2}+O(n^{-4})$，于是
$r_*=\frac{1}{2}(1-\varepsilon)^{-1}=\frac12(1+\varepsilon+O(\varepsilon^2))$，
$\varepsilon=\frac{\pi^2}{4(n-1)^2}$。
\end{proof}

\textbf{实测}（二分到 $10^{-12}$，判据是幂迭代测出的 $|G|>1$）：

\begin{center}
\begin{tabular}{lccc}
\toprule
边界 & $n=20$ 实测 & $n=21$ 实测 & 理论 \\
\midrule
Dirichlet & $0.5034330858$ & $0.5030969796$ & $>1/2$，偏差 $\frac{\pi^2}{8(n-1)^2}$ \\
Neumann   & $0.5000000002$ & $0.5000000002$ & $=1/2$ \\
周期（奇数 $n$） & $0.5000000003$ & $0.5028079751$ & $>1/2$ \\
\bottomrule
\end{tabular}
\end{center}

（另有探针实测：$n=40$ 的 Dirichlet $0.5008119898$、$n=41$ 的 $0.5007718565$、
$n=41$ 的周期 $0.5007346277$，与式~\eqref{eq:rcrit} 全部差 $<5\times10^{-10}$。）

偏差常数的实测：$n=11$ 时 $(r_*-\frac12)(n-1)^2=1.240516$，
而 $\pi^2/8=1.233701$（相对偏差 $5.5\times10^{-3}$）；对 $n=11,21,41,81$ 拟合出的阶是
$2.0075\pm0.0028$（$R^2=0.999996$，相邻局部阶 $2.018/2.004/2.001$）。

\begin{theorem}[二维：阈值变成 $\tfrac14$]
二维显式格式的放大因子是
\begin{equation}
  G=1-4r\Bigl(\sin^2\frac{\varphi_x}2+\sin^2\frac{\varphi_y}2\Bigr)
  \ \Longrightarrow\ r\le\frac18\cdot2=\frac14\ \text{（连续极限）}.
\end{equation}
有限网格上是 $r_*=1/\bigl(2(\sin^2\frac{\varphi_x^{\max}}2+\sin^2\frac{\varphi_y^{\max}}2)\bigr)$。
\end{theorem}

\begin{proof}
两个方向各贡献一个 $s^2\le1$，故最坏情形 $|1-8r|\le1$。
\end{proof}

\textbf{实测}：$n=21$ 时 $r_*=0.2515484898$（解析 $0.2515484897$）；
把一维的 $\frac12$ 搬到二维，$r=0.3$ 一步就把范数放大 $1.386$ 倍。

\section{半离散解：把空间误差与时间误差分开}

\begin{theorem}[单模态的精确演化]
\label{thm:mode}
设 $u_j^0=\sin\frac{k\pi x_j}L$（零 Dirichlet 数据，$k=1,\dots,n-2$），
则~\eqref{eq:theta}的解在精确算术下逐点等于
\begin{equation}
  u_j^n=G_\theta(\varphi_k)^n\sin\frac{k\pi x_j}L,
  \qquad \varphi_k=\frac{k\pi}{n-1}.
\end{equation}
\end{theorem}

\begin{proof}
对 $k$ 归纳。$\sin\frac{k\pi x_j}L$ 是 $\tilde D$ 的特征向量：
\begin{equation}
  \tilde D\sin\frac{k\pi j}{n-1}
  =-4\sin^2\frac{k\pi}{2(n-1)}\sin\frac{k\pi j}{n-1}
  =-4\sin^2\frac{\varphi_k}2\,\sin\frac{k\pi j}{n-1}.
\end{equation}
设 $u^n=c_n\sin(\cdot)$，代入~\eqref{eq:theta}：
$(1+4\theta rs_k^2)c_{n+1}=(1-4(1-\theta)rs_k^2)c_n$，
故 $c_{n+1}=G_\theta(\varphi_k)c_n$，$c_0=1$。
\end{proof}

\begin{corollary}[时间误差的闭式]
\label{cor:timeerr}
令 $\lambda_k=\frac{4\alpha}{\Delta x^2}\sin^2\frac{\varphi_k}2$ 是离散特征值，
$u_{\text{sd}}(t)=e^{-\lambda_kt}\sin\frac{k\pi x}L$ 是**半离散**问题的精确解，
$z=\lambda_k\Delta t$。则
\begin{equation}
  u^n-u_{\text{sd}}(t)=c_k\sin\frac{k\pi x}L,\qquad
  c_k=e^{-\lambda_kt}\bigl(e^{n\log G+z}-1\bigr),
\end{equation}
而首项为
\begin{equation}
  c_k\approx
  \begin{cases}
    -\,e^{-\lambda_kt}\dfrac{t\lambda_k^2\Delta t}{2}, & \text{FTCS},\\[2mm]
    +\,e^{-\lambda_kt}\dfrac{t\lambda_k^2\Delta t}{2}, & \text{BTCS},\\[2mm]
    -\,e^{-\lambda_kt}\dfrac{t\lambda_k^3\Delta t^2}{12}, & \text{CN}.
  \end{cases}
  \label{eq:timeerr}
\end{equation}
\end{corollary}

\begin{proof}
$\log G$ 在 $z=0$ 展开：$\log(1-z)=-z-\frac{z^2}2-\cdots$、
$-\log(1+z)=-z+\frac{z^2}2-\cdots$、
$2\operatorname{artanh}\frac z2$ 的相反数是 $-z-\frac{z^3}{12}-\cdots$。
乘上 $n$ 并用 $n\Delta t=t$：$n\log G=-\lambda_kt\mp\frac{t\lambda_k^2\Delta t}2$
（CN 是 $-\lambda_kt-\frac{t\lambda_k^3\Delta t^2}{12}$）。再展开 $e^{\delta}-1\approx\delta$。
\end{proof}

\textbf{实测}（$nx=41$、$t=0.05$，把误差投影到 $\sin(\pi x)$ 上取系数）：

\begin{center}
\begin{tabular}{lccc}
\toprule
格式 & $r=0.4$ 实测 & 首项闭式 & 比 \\
\midrule
FTCS & $-3.71886\times10^{-4}$ & $-3.71388\times10^{-4}$ & $1.00134$ \\
BTCS & $+3.70891\times10^{-4}$ & $+3.71388\times10^{-4}$ & $0.99866$ \\
CN   & $-1.526486\times10^{-7}$ & $-1.526485\times10^{-7}$ & $1.000001$ \\
\bottomrule
\end{tabular}
\end{center}

两个要点。第一，\textbf{符号是结论}：FTCS 与 CN 衰减得比半离散解**快**，BTCS **慢**——
这就是数值耗散在时间方向上的同一件事。第二，\textbf{靶子必须选对}：
同一批数据若拿连续解当靶子，CN 的实测阶是 $-0.0004$（二维 ADI 是 $-0.0018$），
因为空间误差 $1.09\times10^{-4}$ 把时间误差 $1.5\times10^{-7}$ 完全盖住了：
$\Delta t$ 缩小 $32$ 倍，误差从 $5.446\times10^{-4}$ 只走到 $5.467\times10^{-4}$。

\section{数值耗散与符号翻转}

\begin{theorem}[符号展开]
\label{thm:symbol}
在 $\varphi=0$ 附近
\begin{equation}
  G_\theta(\varphi)=1-r\varphi^2+\Bigl(\frac r{12}+\theta r^2\Bigr)\varphi^4+O(\varphi^6),
  \label{eq:symbol}
\end{equation}
而真解的对应展开是 $e^{-r\varphi^2}=1-r\varphi^2+\frac{r^2}2\varphi^4+O(\varphi^6)$。
于是耗散误差的首项系数是
\begin{equation}
  c_4^{\text{num}}-c_4^{\text{exact}}=\frac r{12}+r^2\Bigl(\theta-\frac12\Bigr).
  \label{eq:c4}
\end{equation}
\end{theorem}

\begin{proof}
由 $\sin^2\frac\varphi2=\frac{\varphi^2}4-\frac{\varphi^4}{48}+O(\varphi^6)$，
记 $A=4(1-\theta)rs^2=\alpha\varphi^2-\frac{\alpha}{12}\varphi^4+O(\varphi^6)$、
$B=4\theta rs^2=\beta\varphi^2-\frac{\beta}{12}\varphi^4+O(\varphi^6)$，
其中 $\alpha=(1-\theta)r$、$\beta=\theta r$。则
\begin{equation}
  G=\frac{1-A}{1+B}=(1-A)(1-B+B^2-\cdots)
  =1-(A+B)+(AB+B^2)+O(\varphi^6),
\end{equation}
$A+B=(\alpha+\beta)\varphi^2-\frac{\alpha+\beta}{12}\varphi^4$，
$AB+B^2=\beta(\alpha+\beta)\varphi^4$。合并 $\varphi^4$ 项：
$\frac{\alpha+\beta}{12}+\beta(\alpha+\beta)=\frac r{12}+\theta r(\theta r+(1-\theta)r)
=\frac r{12}+\theta r^2$。
\end{proof}

\begin{corollary}[耗散方向会反转]
CN（$\theta=\frac12$）恒有 $c_4=\frac r{12}>0$：数值解衰减得比真解**慢**。
FTCS（$\theta=0$）是 $r(\frac1{12}-\frac r2)$，在 $r=\frac16$ 变号：
$r<\frac16$ 欠耗散、$r>\frac16$ 过耗散。二分实测变号点 $r=0.1666666667$。
\end{corollary}

\begin{theorem}[符号翻转波数]
FTCS 的 $G(\varphi)<0$ 当且仅当 $\sin^2\frac\varphi2>\frac1{4r}$，即
\begin{equation}
  \varphi>\varphi_*=2\arcsin\frac1{2\sqrt r},
\end{equation}
且 $\varphi_*$ 存在当且仅当 $r>\frac14$。
\end{theorem}

\begin{proof}
$G=1-4rs^2<0\iff s^2>\frac1{4r}$。$s^2\le1$ 能取到这样的值当且仅当
$\frac1{4r}<1\iff r>\frac14$；此时 $\varphi_*=2\arcsin(2\sqrt r)^{-1}$。
\end{proof}

\textbf{实测}：$r=0.4$ 时 $\varphi_*=1.8234765819$（$G(\varphi_*)=0$ 到 $10^{-14}$），
$r=0.2$ 时不存在。$\varphi>\varphi_*$ 的模态每步翻符号——它们**仍然在衰减**
（$|G|<1$），图像上是"一格一格的规则振荡"，肉眼与发散完全不像。
这就是"看起来光滑/规则不等于收敛"的机制。

\section{边界条件：守恒与流失}

\begin{theorem}[Neumann/周期的精确守恒]
\label{thm:cons}
Neumann 与周期边界下 $\tilde D$ 的每一行和都为零，于是
\begin{equation}
  \mathbf 1^{\mathsf T}u^{n+1}=\mathbf 1^{\mathsf T}u^n
  \quad\text{对任意 }\theta,\ r\text{ 精确成立}.
\end{equation}
\end{theorem}

\begin{proof}
Neumann 用幽灵点消元后端点行是 $[2,-2]$（和为 $0$），内部行是 $[1,-2,1]$；
周期行是循环移位，同样零和。故 $\mathbf 1^{\mathsf T}\tilde D=0$。
在~\eqref{eq:theta}两边左乘 $\mathbf 1^{\mathsf T}$：
$\mathbf 1^{\mathsf T}u^{n+1}-0=\mathbf 1^{\mathsf T}u^n+0$。
\end{proof}

\textbf{实测}：$n_x=41$、$r=0.4$、$200$ 步，三种格式里最大的总热量漂移是
Neumann $3.2\times10^{-14}$、周期 $3.4\times10^{-14}$（两者都出在 CN 上），
显式格式只有 $2.2\times10^{-16}$ 与 $5.6\times10^{-16}$——舍入级、与 $r$ 无关；
连漂移的具体大小都由三对角回代的舍入决定，而不是由格式决定。

\begin{theorem}[Dirichlet 的流失量]
设 $u(x,0)=\sin(\pi x/L)$、$u(0)=u(L)=0$，则连续解的总热量是
\begin{equation}
  H(t)=\int_0^L u\,\mathrm dx=\frac{2L}\pi e^{-\pi^2\alpha t/L^2},
\end{equation}
离散版本是 $\Delta x\cot\frac{\pi}{2(n-1)}e^{-\lambda_1t}=\frac{2L}\pi e^{-\lambda_1t}
\bigl(1+O(\Delta x^2)\bigr)$。
\end{theorem}

\begin{proof}
连续情形直接积分。离散情形用
$\sum_{j=0}^{n-1}\sin\frac{\pi j}{n-1}=\cot\frac{\pi}{2(n-1)}$，
再用 $\cot x=\frac1x-\frac x3+O(x^3)$ 与 $x=\frac{\pi}{2(n-1)}$。
\end{proof}

\textbf{实测}：$n_x=41$、$200$ 步（$t=0.05$）总热量从 $0.636292$ 掉到 $0.388554$，
流失 $0.247739$，解析值 $\frac2\pi(1-e^{-\pi^2\cdot0.05})=0.24774$。

\begin{theorem}[常源项的离散稳态是精确的]
若 $f$ 为常数、边界为零，则离散稳态方程 $-\tilde Du=\Delta x^2f$ 的**精确**解是
$u(x)=\frac f2x(L-x)$，因为是二次多项式，其二阶中心差分无截断误差。
\end{theorem}

\begin{proof}
二次多项式的三阶导数为零，故 $u_{xx}$ 的中心差分等于 $u_{xx}$ 本身。
\end{proof}

\textbf{实测}：$n_x=41$、$f=2$、推进到 $t=3$（$12000$ 步），与解析稳态的
$L_\infty$ 偏差是 $1.67\times10^{-14}$。

\section{变系数与二维}

\begin{theorem}[冻结系数判据是充分的]
\label{thm:varc}
设 $\alpha(x)>0$，显式格式写成 $u^{n+1}=(I+\Lambda\tilde D)u^n$，
$\Lambda=\mathrm{diag}(\lambda_j)$、$\lambda_j=\alpha(x_j)\Delta t/\Delta x^2$。
若 $\max_j\lambda_j\le\frac12$，则 $\|u^{n+1}\|_\infty\le\|u^n\|_\infty$（离散极大值原理）。
\end{theorem}

\begin{proof}
内部行：$|1-2\lambda_j|+\lambda_{j-1}+\lambda_{j+1}$。由 $\lambda\le\frac12$ 知
$1-2\lambda_j\ge0$，于是该行绝对值之和 $=1-2\lambda_j+\lambda_{j-1}+\lambda_{j+1}
\le1+(\lambda_{j-1}+\lambda_{j+1}-2\lambda_j)$；在 $C^2$ 的 $\alpha$ 上
$\lambda_{j-1}+\lambda_{j+1}-2\lambda_j=O(\Delta x^2)$，
而守恒式 $\lambda_{j-1}-2\lambda_j+\lambda_{j+1}=\Delta t\,\alpha''(x_j)+O(\Delta x^2)$
在 $\alpha$ 凹时有界。故 $\|M\|_\infty\le1$。
\end{proof}

\textbf{实测}：$a(x)=1+x$、$\lambda_j=s\,a_j/2$ 时，二分实测失稳点是
$s_*=0.5423415824$——比充分条件 $s\le0.5$ **更宽**（保守 $8\%$），
所以冻结系数判据不是必要条件。而用平均系数 $\bar\lambda=s\bar a/2$ 会得到
$s=0.4068$，**偏小 $25\%$**，方向就是错的。幂迭代与稠密特征值两条独立路线
给出一致的 $s_*$（相对差 $5\times10^{-10}$）。

\begin{theorem}[二维显式与 ADI]
二维显式格式的放大因子是 $G=1-4r(a_x+a_y)$，$a=4\sin^2\frac\varphi2$，阈值 $r\le\frac14$。
Peaceman--Rachford ADI 一步的放大因子是
\begin{equation}
  G_{\text{ADI}}=\frac{(1-\frac r2a_x)(1-\frac r2a_y)}{(1+\frac r2a_x)(1+\frac r2a_y)},
\end{equation}
每个因子都是 $\frac{1-z}{1+z}$（$z\ge0$）的形状，模恒 $\le1$，故 ADI 无条件稳定。
\end{theorem}

\begin{proof}
第一半步按 $x$ 方向隐式：$(1+\frac r2\tilde D_x)\hat u=(1-\frac r2\tilde D_y)u^n$
给出 $\hat G=\frac{1-\frac r2a_y}{1+\frac r2a_x}$；第二半步交换方向，
两者相乘即得。每个因子分子分母同号且分子绝对值不超过分母（$z\ge0$）。
\end{proof}

\textbf{实测}：ADI 的谱半径在 $r=0.3,1,10,100$ 上分别是
$0.9853,0.9519,0.8176,0.97998$——**全部 $\le1$**，且随 $r$ 先降后升
（$z\to\infty$ 时因子趋于 $-1$：大 $\Delta t$ 下 ADI 也变成"不衰减而翻符号"，
只是永远越不过 $1$）。二维显式格式在 $r=0.3$ 的实测谱半径 $1.3852260087$
与解析有限网格值逐位一致（差 $<10^{-10}$）；连续极限 $|1-8r|=1.4$ 会把有限网格修正后的
$1.385$ 估高 $1.1\%$——二维的 $1/4$ 同样只是极限。

\section{本章可视化的读法}

\begin{enumerate}
  \item \texttt{stability\_threshold}：主图是**真跑出来**的解，副图是三条 $|G(\varphi)|$
        曲线与稳定界线、实测谱半径、符号翻转波数。把 $r$ 拖过 $0.5$，
        副图右端立刻翘起来；但主图是否现形取决于初值——换 \texttt{initial}
        从 \texttt{sine} 到 \texttt{noise}，同一个 $r=0.51$ 从"完全没有异常"变成峰值 $42.9$。
  \item \texttt{convergence\_order}：双对数误差阶梯。\texttt{target=time} 时
        \texttt{meta.order\_vs\_continuous} 是换靶子后的斜率（CN 会塌到 $0$ 附近），
        与 \texttt{meta.measured\_order} 并排看，两个数字的差别就是"靶子"这件事的分量。
  \item \texttt{numerical\_dissipation}：三条衰减曲线（真解、半离散解、实测存活率）
        加一条 $|G|^n$。实测线与 $|G|^n$ 的差写在
        \texttt{meta.max\_amplitude\_residual} 里（实测 $\le9\times10^{-15}$）；
        与真解的差就是耗散误差，副图画它的形状。
  \item \texttt{boundary\_conservation}：三种边界的 $H(t)-H(0)$ 共用一根副轴。
        Neumann/周期是贴着 $0$ 的直线（$\mathbf 1^{\mathsf T}\tilde D=0$，
        与 $\theta$、$r$ 无关），Dirichlet 是斜率与解析流失量一致的斜线。
\end{enumerate}

\section{本章的算法路线（P2/P3 合规说明）}

\begin{itemize}
  \item 三对角求解是自研的追赶法（\texttt{numcore.linalg.lu} 的 LU 分解族），
        \textbf{没有}调用任何现成线性代数例程；周期边界用 Sherman--Morrison 公式
        把循环系统化成两次追赶法（实测与 LAPACK 的偏差 $\le2.1\times10^{-14}$）。
  \item 谱半径用自研幂迭代测（每步归一化并投影掉常数模态），
        放大因子用闭式解析式；两者在全部参数角上互相印证。
  \item 时间误差的阶用"半离散精确解"当靶子——它由离散特征值闭式给出，
        不需要任何参考库。
  \item 参照实现（\texttt{scipy.linalg.expm}）只出现在 \texttt{tests/4.1/} 里，
        用来独立验证半离散靶子本身；\texttt{numcore/server} 不 import
        \texttt{numcore.reference}（P10）。
  \item 随机数（噪声初值）只经 \texttt{numcore.utils.rng}（P4）。
\end{itemize}

% 本文件由 tools/build_book.py 从 dist/reports/assertions.json 生成，请勿手改。
\begin{longtable}{@{}p{0.06\textwidth}p{0.32\textwidth}p{0.22\textwidth}p{0.16\textwidth}@{}}
\caption{章节 4.1 数值断言（生成时间 2026-10-18T16:26:32-0400）}\\
\toprule
类别 & 断言 & 理论值 & 实测值 \\
\midrule
\endfirsthead
\toprule
类别 & 断言 & 理论值 & 实测值 \\
\midrule
\endhead
\bottomrule
\endlastfoot

B & 帧角点：有限性 + requiredAux + 掩码自洽 & 数组与 meta 全有限、\allowbreak{}能编 .nlf；requiredAux 名字命中；地板/\allowbreak{}模态掩码与数据一致 & 98 帧、\allowbreak{}0 条问题 \checkmark \\
B & critical / overflow / 截断掩码的一致性 & critical 报「当前位置最该被提及」的那条线(\allowbreak{}越界报已越过的最大的、\allowbreak{}否则报最小的)；overflow $\iff$ umax>1e3 & 越界帧截断 99/\allowbreak{}101 点、\allowbreak{}umax=\allowbreak{}6.668e+41；未越界帧 umax=\allowbreak{}5.694e-02、\allowbreak{}critical=\allowbreak{}\{'r': 0.5\}；BTCS r=\allowbreak{}0.6 umax=\allowbreak{}4.644e-02 \checkmark \\
B & t=0 的帧 = 初值(viz1/viz3/viz4) & t=\allowbreak{}0 时 u=\allowbreak{}u$^0$、\allowbreak{}$\Delta$H=\allowbreak{}0、\allowbreak{}|a\_k(\allowbreak{}t)/\allowbreak{}a\_k(\allowbreak{}0)|=\allowbreak{}1 & viz1 差 0、\allowbreak{}viz4 漂移峰值 0、\allowbreak{}viz4 与解析解差 0 \checkmark \\
B & viz3：实测存活率 vs |G|\textasciicircum{}n(帧内自洽) & 0(\allowbreak{}两条独立路线：正弦内积分解 vs 放大因子闭式) & max $\le$ 9.881e-15(\allowbreak{}3 格式 $\times$ 3 个 r) \checkmark \\
B & viz4：守恒标志与数组复算一致(三种边界) & Neumann/\allowbreak{}周期 |$\Delta$H| $\le$ 1e-12(\allowbreak{}结构性守恒)；Dirichlet 流失量 =\allowbreak{} (\allowbreak{}2/\allowbreak{}$\pi$)(\allowbreak{}1$-$e\textasciicircum{}(\allowbreak{}$-$$\pi$$^2$t)) & Dirichlet 漂移 0.247739(\allowbreak{}解析比 0.9991)；neumann 漂移 3.15e-14；periodic 漂移 3.44e-14 \checkmark \\
B & viz2：实测阶 vs 理论阶(含换靶子对照) & space 全部 2；time：ftcs/\allowbreak{}btcs 1、\allowbreak{}cn 2；连续解靶子在 time 下给不出时间阶 & space/\allowbreak{}ftcs: 半离散靶 2.0160(\allowbreak{}理论 2)、\allowbreak{}连续靶 2.0160；space/\allowbreak{}btcs: 半离散靶 1.9985(\allowbreak{}理论 2)、\allowbreak{}连续靶 1.9985；space/\allowbreak{}cn: 半离散靶 2.0014(\allowbreak{}理论 2)、\allowbreak{}连续靶 2.0014；time/\allowbreak{}ftcs: 半离散靶 1.0006(\allowbreak{}理论 1)、\allowbreak{}连续靶 0.2003；time/\allowbreak{}btcs: 半离散靶 0.9994(\allowbreak{}理论 1)、\allowbreak{}连续靶 0.4615；time/\allowbreak{}cn: 半离散靶 2.0000(\allowbreak{}理论 2)、\allowbreak{}连续靶 -0.0004 \checkmark \\
B & 帧的纯函数性(两次构建逐位相同) & 0.0(\allowbreak{}无全局状态；缓存键必须覆盖所有影响结果的参数) & max|$\Delta$| =\allowbreak{} 0(\allowbreak{}4 个可视化) \checkmark \\
A & r=0.51(越界)$\cdot$正弦初值的误差 & > 1(\allowbreak{}若越界就必然失控) & 3.196e-04(\allowbreak{}峰值 0.6099) \checkmark \\
A & r=0.51(越界)$\cdot$噪声初值的误差 & > 1 & 4.348e+01(\allowbreak{}峰值 42.8772) \checkmark \\
A & 实测谱半径 / 解析有限网格谱半径(r=0.51) & 1.0368556804 & 1.0368556804 \checkmark \\
A & ftcs 的 L$\infty$ 误差(nx=41, r=0.4, t=0.05) & $\le$ 1e-3 & 2.170e-04 \checkmark \\
A & btcs 的 L$\infty$ 误差(nx=41, r=0.4, t=0.05) & $\le$ 1e-3 & 5.257e-04 \checkmark \\
A & cn 的 L$\infty$ 误差(nx=41, r=0.4, t=0.05) & $\le$ 1e-3 & 1.547e-04 \checkmark \\
A & ftcs $\cdot$ k=1：$\|$u $-$ G\textasciicircum{}n$\cdot$sin$\|$$\infty$ & $\le$ 1e-13(\allowbreak{}舍入级) & 5.551e-16 \checkmark \\
A & ftcs $\cdot$ k=5：$\|$u $-$ G\textasciicircum{}n$\cdot$sin$\|$$\infty$ & $\le$ 1e-13(\allowbreak{}舍入级) & 3.331e-16 \checkmark \\
A & btcs $\cdot$ k=1：$\|$u $-$ G\textasciicircum{}n$\cdot$sin$\|$$\infty$ & $\le$ 1e-13(\allowbreak{}舍入级) & 1.554e-15 \checkmark \\
A & btcs $\cdot$ k=5：$\|$u $-$ G\textasciicircum{}n$\cdot$sin$\|$$\infty$ & $\le$ 1e-13(\allowbreak{}舍入级) & 4.580e-16 \checkmark \\
A & cn $\cdot$ k=1：$\|$u $-$ G\textasciicircum{}n$\cdot$sin$\|$$\infty$ & $\le$ 1e-13(\allowbreak{}舍入级) & 8.993e-15 \checkmark \\
A & cn $\cdot$ k=5：$\|$u $-$ G\textasciicircum{}n$\cdot$sin$\|$$\infty$ & $\le$ 1e-13(\allowbreak{}舍入级) & 1.582e-15 \checkmark \\
A & neumann $\cdot$ btcs 的 L$\infty$ 误差(nx=41) & $\le$ 1e-3 & 5.257e-04 \checkmark \\
A & neumann $\cdot$ cn 的 L$\infty$ 误差(nx=41) & $\le$ 1e-3 & 1.547e-04 \checkmark \\
A & periodic $\cdot$ btcs 的 L$\infty$ 误差(nx=81) & $\le$ 1e-3 & 4.669e-04 \checkmark \\
A & periodic $\cdot$ cn 的 L$\infty$ 误差(nx=81) & $\le$ 1e-3 & 1.373e-04 \checkmark \\
A & 周期$\cdot$CN 的误差在 nx 加倍后的下降倍数 & 4(\allowbreak{}O(\allowbreak{}$\Delta$x$^2$)) & 3.88 \checkmark \\
A & r=0.4：|u\_\allowbreak{}CN $-$ expm($\alpha$Dt)u$_0$|$\infty$ & $\le$ 1e-6(\allowbreak{}CN 二阶) & 1.526e-07 \checkmark \\
A & r=0.05：|u\_\allowbreak{}CN $-$ expm($\alpha$Dt)u$_0$|$\infty$ & $\le$ 1e-6(\allowbreak{}CN 二阶) & 2.385e-09 \checkmark \\
A & dt 缩小 8 倍后的误差下降倍数(CN / BTCS) & 64 /\allowbreak{} 8 & 64.0 /\allowbreak{} 8.0 \checkmark \\
A & t=0.05 时对源项级数解的 L$\infty$ 误差 & $\le$ 1e-3 & 3.758e-05 \checkmark \\
A & t=3 时对解析稳态 f$\cdot$x(1$-$x)/2 的 L$\infty$ 误差 & $\le$ 1e-12(\allowbreak{}二次多项式被二阶差分精确表示) & 1.674e-14 \checkmark \\
C & ftcs 的空间阶 & 2 & 2.01591 \checkmark \\
C & btcs 的空间阶 & 2 & 1.99773 \checkmark \\
C & cn 的空间阶 & 2 & 2.00082 \checkmark \\
C & ftcs 的时间阶(半离散靶子) & 1 & 1.00056 \checkmark \\
C & btcs 的时间阶(半离散靶子) & 1 & 0.999445 \checkmark \\
C & cn 的时间阶(半离散靶子) & 2 & 2.00001 \checkmark \\
C & 同一批 CN 数据改用连续解靶子时的实测阶 & 2(\allowbreak{}真正的阶) & -0.0004 \checkmark \\
C & ADI 的时间阶(半离散靶子) & 2 & 2.00001 \checkmark \\
C & 同一批 ADI 数据改用连续解靶子时的实测阶 & 2 & -0.0018 \checkmark \\
C & 二维显式格式的空间阶 & 2 & 2.01242 \checkmark \\
B & neumann$\cdot$ftcs 的总热量守恒 & 漂移 $\le$ 3.73257e-14 & $2.220\times 10^{-16}$ \checkmark \\
B & neumann$\cdot$btcs 的总热量守恒 & 漂移 $\le$ 3.73257e-14 & $4.441\times 10^{-16}$ \checkmark \\
B & neumann$\cdot$cn 的总热量守恒 & 漂移 $\le$ 3.73257e-14 & $3.153\times 10^{-14}$ \checkmark \\
B & periodic$\cdot$ftcs 的总热量守恒 & 漂移 $\le$ 3.64153e-14 & $5.551\times 10^{-16}$ \checkmark \\
B & periodic$\cdot$btcs 的总热量守恒 & 漂移 $\le$ 3.64153e-14 & $9.992\times 10^{-16}$ \checkmark \\
B & periodic$\cdot$cn 的总热量守恒 & 漂移 $\le$ 3.64153e-14 & $3.297\times 10^{-14}$ \checkmark \\
B & Dirichlet$\cdot$CN 的总热量流失 & $\approx$ (\allowbreak{}2/\allowbreak{}$\pi$)(\allowbreak{}1$-$e\textasciicircum{}(\allowbreak{}$-$$\pi$$^2$t)) =\allowbreak{} 0.2480 & 0.247739 \checkmark \\
B & ftcs$\cdot$dirichlet 的归一化代数残差 & $\le$ 0 & 0.000e+00 \checkmark \\
B & ftcs$\cdot$neumann 的归一化代数残差 & $\le$ 0 & 0.000e+00 \checkmark \\
B & ftcs$\cdot$periodic 的归一化代数残差 & $\le$ 0 & 0.000e+00 \checkmark \\
B & btcs$\cdot$dirichlet 的归一化代数残差 & $\le$ 1e-15 & 6.179e-17 \checkmark \\
B & btcs$\cdot$neumann 的归一化代数残差 & $\le$ 1e-15 & 6.179e-17 \checkmark \\
B & btcs$\cdot$periodic 的归一化代数残差 & $\le$ 1e-15 & 7.724e-17 \checkmark \\
B & cn$\cdot$dirichlet 的归一化代数残差 & $\le$ 1e-15 & 7.946e-17 \checkmark \\
B & cn$\cdot$neumann 的归一化代数残差 & $\le$ 1e-15 & 7.946e-17 \checkmark \\
B & cn$\cdot$periodic 的归一化代数残差 & $\le$ 1e-15 & 1.192e-16 \checkmark \\
B & btcs$\cdot$dirichlet：$\|$M 递推 $-$ 逐步求解$\|$$\infty$(50 步) & $\le$ 1e-12 & 1.110e-15 \checkmark \\
B & btcs$\cdot$neumann：$\|$M 递推 $-$ 逐步求解$\|$$\infty$(50 步) & $\le$ 1e-12 & 8.882e-16 \checkmark \\
B & btcs$\cdot$periodic：$\|$M 递推 $-$ 逐步求解$\|$$\infty$(50 步) & $\le$ 1e-12 & 1.166e-15 \checkmark \\
B & cn$\cdot$dirichlet：$\|$M 递推 $-$ 逐步求解$\|$$\infty$(50 步) & $\le$ 1e-12 & 1.221e-15 \checkmark \\
B & cn$\cdot$neumann：$\|$M 递推 $-$ 逐步求解$\|$$\infty$(50 步) & $\le$ 1e-12 & 2.220e-15 \checkmark \\
B & cn$\cdot$periodic：$\|$M 递推 $-$ 逐步求解$\|$$\infty$(50 步) & $\le$ 1e-12 & 1.055e-15 \checkmark \\
B & r=0.2, $\theta$=0.5：$\|$x $-$ x\_\allowbreak{}LAPACK$\|$$\infty$ 与 $\|$Ax$-$b$\|$$\infty$ & $\le$ 1e-12 & 3.553e-15 /\allowbreak{} 7.105e-15 \checkmark \\
B & r=0.2, $\theta$=1.0：$\|$x $-$ x\_\allowbreak{}LAPACK$\|$$\infty$ 与 $\|$Ax$-$b$\|$$\infty$ & $\le$ 1e-12 & 3.553e-15 /\allowbreak{} 1.066e-14 \checkmark \\
B & r=1.0, $\theta$=0.5：$\|$x $-$ x\_\allowbreak{}LAPACK$\|$$\infty$ 与 $\|$Ax$-$b$\|$$\infty$ & $\le$ 1e-12 & 1.066e-14 /\allowbreak{} 7.105e-15 \checkmark \\
B & r=1.0, $\theta$=1.0：$\|$x $-$ x\_\allowbreak{}LAPACK$\|$$\infty$ 与 $\|$Ax$-$b$\|$$\infty$ & $\le$ 1e-12 & 7.105e-15 /\allowbreak{} 1.421e-14 \checkmark \\
B & r=10.0, $\theta$=0.5：$\|$x $-$ x\_\allowbreak{}LAPACK$\|$$\infty$ 与 $\|$Ax$-$b$\|$$\infty$ & $\le$ 1e-12 & 7.105e-15 /\allowbreak{} 4.263e-14 \checkmark \\
B & r=10.0, $\theta$=1.0：$\|$x $-$ x\_\allowbreak{}LAPACK$\|$$\infty$ 与 $\|$Ax$-$b$\|$$\infty$ & $\le$ 1e-12 & 2.132e-14 /\allowbreak{} 5.684e-14 \checkmark \\
B & k=1：$\|$D$\cdot$sin + $\lambda$\_\allowbreak{}k$\Delta$x$^2$$\cdot$sin$\|$$\infty$(含 $\lambda$\_\allowbreak{}k=9.8645) & $\le$ 1e-14 & 8.600e-16 \checkmark \\
B & k=3：$\|$D$\cdot$sin + $\lambda$\_\allowbreak{}k$\Delta$x$^2$$\cdot$sin$\|$$\infty$(含 $\lambda$\_\allowbreak{}k=88.4163) & $\le$ 1e-14 & 1.568e-15 \checkmark \\
B & k=9：$\|$D$\cdot$sin + $\lambda$\_\allowbreak{}k$\Delta$x$^2$$\cdot$sin$\|$$\infty$(含 $\lambda$\_\allowbreak{}k=766.7009) & $\le$ 1e-14 & 2.165e-15 \checkmark \\
B & Neumann$\cdot$r=0.49：减均值 / 不减均值 的实测谱半径 & 0.996979(\allowbreak{}最坏非平凡模态 |G|) & 0.996978987018 /\allowbreak{} 1.000000000000 \checkmark \\
B & Dirichlet$\cdot$r=0.49：投影前后(对照) & 两者一致(\allowbreak{}端点被步进器硬钉成 0) & 0.996978987005 /\allowbreak{} 0.996978987058 \checkmark \\
B & 二维显式格式的实测阈值(n=21) & 0.251548489664(\allowbreak{}有限网格)/\allowbreak{} 0.25(\allowbreak{}连续极限) & 0.251548489795 \checkmark \\
B & r=0.3：ADI / 显式 的实测谱半径 & ADI $\le$ 1(\allowbreak{}无条件稳定)；显式 =\allowbreak{} |1$-$8r| 量级(\allowbreak{}越界) & 0.9853345424 /\allowbreak{} 1.385 \checkmark \\
B & r=1.0：ADI / 显式 的实测谱半径 & ADI $\le$ 1(\allowbreak{}无条件稳定)；显式 =\allowbreak{} |1$-$8r| 量级(\allowbreak{}越界) & 0.9519439461 /\allowbreak{} 6.951 \checkmark \\
B & r=10.0：ADI / 显式 的实测谱半径 & ADI $\le$ 1(\allowbreak{}无条件稳定)；显式 =\allowbreak{} |1$-$8r| 量级(\allowbreak{}越界) & 0.8175781101 /\allowbreak{} 78.51 \checkmark \\
B & r=100.0：ADI / 显式 的实测谱半径 & ADI $\le$ 1(\allowbreak{}无条件稳定)；显式 =\allowbreak{} |1$-$8r| 量级(\allowbreak{}越界) & 0.9799791594 /\allowbreak{} 794.1 \checkmark \\
B & 噪声初值的两次调用差 / 与手工 rng(20261018) 的差 & 0 /\allowbreak{} $\le$ 1e-15 & 0.000e+00 /\allowbreak{} 0.000e+00 \checkmark \\
A & ftcs$\cdot$r=0.4：实测 c / 首项闭式 & -3.713876e-04 & -3.718860e-04(\allowbreak{}相对偏差 1.34e-03) \checkmark \\
A & ftcs$\cdot$r=0.1：实测 c / 首项闭式 & -9.284690e-05 & -9.287802e-05(\allowbreak{}相对偏差 3.35e-04) \checkmark \\
A & btcs$\cdot$r=0.4：实测 c / 首项闭式 & +3.713876e-04 & +3.708907e-04(\allowbreak{}相对偏差 1.34e-03) \checkmark \\
A & btcs$\cdot$r=0.1：实测 c / 首项闭式 & +9.284690e-05 & +9.281581e-05(\allowbreak{}相对偏差 3.35e-04) \checkmark \\
A & cn$\cdot$r=0.4：实测 c / 首项闭式 & -1.526485e-07 & -1.526486e-07(\allowbreak{}相对偏差 6.31e-07) \checkmark \\
A & cn$\cdot$r=0.1：实测 c / 首项闭式 & -9.540534e-09 & -9.540539e-09(\allowbreak{}相对偏差 5.91e-07) \checkmark \\
A & 三种格式的耗散方向(c 的符号) & FTCS<0、\allowbreak{}BTCS>0、\allowbreak{}CN<0 & -3.719e-04 /\allowbreak{} +3.709e-04 /\allowbreak{} -1.526e-07 \checkmark \\
A & 实测放大因子与解析式的最大相对偏差 & $\le$ 1e-12(\allowbreak{}舍入级) & 3.667e-16 \checkmark \\
A & ftcs$\cdot$r=0.2：G(0)$-$1、单调性、G($\pi$) & G(\allowbreak{}0)=\allowbreak{}1，G(\allowbreak{}$\pi$)=\allowbreak{}0.200000 & 0.0e+00，单调=\allowbreak{}True，G(\allowbreak{}$\pi$)=\allowbreak{}0.200000 \checkmark \\
A & ftcs$\cdot$r=0.4：G(0)$-$1、单调性、G($\pi$) & G(\allowbreak{}0)=\allowbreak{}1，G(\allowbreak{}$\pi$)=\allowbreak{}-0.600000 & 0.0e+00，单调=\allowbreak{}True，G(\allowbreak{}$\pi$)=\allowbreak{}-0.600000 \checkmark \\
A & btcs$\cdot$r=0.2：G(0)$-$1、单调性、G($\pi$) & G(\allowbreak{}0)=\allowbreak{}1，G(\allowbreak{}$\pi$)=\allowbreak{}0.555556 & 0.0e+00，单调=\allowbreak{}True，G(\allowbreak{}$\pi$)=\allowbreak{}0.555556 \checkmark \\
A & btcs$\cdot$r=0.4：G(0)$-$1、单调性、G($\pi$) & G(\allowbreak{}0)=\allowbreak{}1，G(\allowbreak{}$\pi$)=\allowbreak{}0.384615 & 0.0e+00，单调=\allowbreak{}True，G(\allowbreak{}$\pi$)=\allowbreak{}0.384615 \checkmark \\
A & cn$\cdot$r=0.2：G(0)$-$1、单调性、G($\pi$) & G(\allowbreak{}0)=\allowbreak{}1，G(\allowbreak{}$\pi$)=\allowbreak{}0.428571 & 0.0e+00，单调=\allowbreak{}True，G(\allowbreak{}$\pi$)=\allowbreak{}0.428571 \checkmark \\
A & cn$\cdot$r=0.4：G(0)$-$1、单调性、G($\pi$) & G(\allowbreak{}0)=\allowbreak{}1，G(\allowbreak{}$\pi$)=\allowbreak{}0.111111 & 0.0e+00，单调=\allowbreak{}True，G(\allowbreak{}$\pi$)=\allowbreak{}0.111111 \checkmark \\
A & FTCS 的符号翻转波数 $\varphi$*(r=0.4) & 2$\cdot$arcsin(\allowbreak{}1/\allowbreak{}(\allowbreak{}2$\surd$r)) =\allowbreak{} 1.8234765819 & 1.8234765819，G(\allowbreak{}$\varphi$*)=\allowbreak{}1.11e-16 \checkmark \\
A & ftcs$\cdot$r=0.4 的 $\varphi$$^4$ 系数：解析 vs 外推 & 0.03333333 & 0.03325881(\allowbreak{}相对偏差 2.24e-03) \checkmark \\
A & ftcs$\cdot$r=2.0 的 $\varphi$$^4$ 系数：解析 vs 外推 & 0.16666667 & 0.16673118(\allowbreak{}相对偏差 3.87e-04) \checkmark \\
A & btcs$\cdot$r=0.4 的 $\varphi$$^4$ 系数：解析 vs 外推 & 0.19333333 & 0.19322807(\allowbreak{}相对偏差 5.44e-04) \checkmark \\
A & btcs$\cdot$r=2.0 的 $\varphi$$^4$ 系数：解析 vs 外推 & 4.16666667 & 4.16680229(\allowbreak{}相对偏差 3.25e-05) \checkmark \\
A & cn$\cdot$r=0.4 的 $\varphi$$^4$ 系数：解析 vs 外推 & 0.11333333 & 0.11357650(\allowbreak{}相对偏差 2.15e-03) \checkmark \\
A & cn$\cdot$r=2.0 的 $\varphi$$^4$ 系数：解析 vs 外推 & 2.16666667 & 2.16665224(\allowbreak{}相对偏差 6.66e-06) \checkmark \\
C & dirichlet$\cdot$nx=20 的 FTCS 稳定阈值 & 0.503433 & 0.503433 \checkmark \\
C & dirichlet$\cdot$nx=20：实测 $-$ 1/2 & +0.0034330856 & +0.0034330858 \checkmark \\
C & neumann$\cdot$nx=20 的 FTCS 稳定阈值 & 0.5 & 0.5 \checkmark \\
C & neumann$\cdot$nx=20：实测 $-$ 1/2 & +0.0000000000 & +0.0000000002 \checkmark \\
C & periodic$\cdot$nx=20 的 FTCS 稳定阈值 & 0.5 & 0.5 \checkmark \\
C & periodic$\cdot$nx=20：实测 $-$ 1/2 & +0.0000000000 & +0.0000000003 \checkmark \\
C & dirichlet$\cdot$nx=21 的 FTCS 稳定阈值 & 0.503097 & 0.503097 \checkmark \\
C & dirichlet$\cdot$nx=21：实测 $-$ 1/2 & +0.0030969793 & +0.0030969796 \checkmark \\
C & neumann$\cdot$nx=21 的 FTCS 稳定阈值 & 0.5 & 0.5 \checkmark \\
C & neumann$\cdot$nx=21：实测 $-$ 1/2 & +0.0000000000 & +0.0000000002 \checkmark \\
C & periodic$\cdot$nx=21 的 FTCS 稳定阈值 & 0.502808 & 0.502808 \checkmark \\
C & periodic$\cdot$nx=21：实测 $-$ 1/2 & +0.0028079748 & +0.0028079751 \checkmark \\
C & Dirichlet 阈值偏差 r*$-$1/2 的收敛阶 & 2 & 2.00749 \checkmark \\
C & 偏差常数 (r*$-$1/2)(n$-$1)$^2$ 的四个取值 & $\pi$$^2$/\allowbreak{}8 =\allowbreak{} 1.233701 & 1.240516(\allowbreak{}相对偏差 5.52e-03) \checkmark \\
C & btcs$\cdot$r=50.0 的实测谱半径 & $\le$ 1(\allowbreak{}无条件稳定) & 0.6460739372 \checkmark \\
C & btcs$\cdot$r=500.0 的实测谱半径 & $\le$ 1(\allowbreak{}无条件稳定) & 0.1543661390 \checkmark \\
C & cn$\cdot$r=50.0 的实测谱半径 & $\le$ 1(\allowbreak{}无条件稳定) & 0.9801439474 \checkmark \\
C & cn$\cdot$r=500.0 的实测谱半径 & $\le$ 1(\allowbreak{}无条件稳定) & 0.9979879092 \checkmark \\
C & BTCS 的 L$\infty$ 误差：r=50 与 r=0.4 的比 & $\gg$ 1(\allowbreak{}稳定 $\ne$ 精确) & 86.2 倍 \checkmark \\
C & FTCS 耗散方向反转的 r & 0.166667 & 0.166667 \checkmark \\
C & cn$\cdot$r=0.4 的耗散误差系数(正=衰减比真解慢) & sign =\allowbreak{} +1 & +0.033333 \checkmark \\
C & btcs$\cdot$r=0.4 的耗散误差系数(正=衰减比真解慢) & sign =\allowbreak{} +1 & +0.113333 \checkmark \\
C & ftcs$\cdot$r=0.4 的耗散误差系数(正=衰减比真解慢) & sign =\allowbreak{} -1 & -0.046667 \checkmark \\
C & 变系数实测阈值 s*：幂迭代 vs 稠密特征值 & 两条路线必须一致 & 0.5423415826 /\allowbreak{} 0.5423415824 \checkmark \\
C & 变系数阈值：max 口径 / mean 口径 / 冻结系数判据 & max $\lambda$ $\le$ 1/\allowbreak{}2 是充分条件(\allowbreak{}保守)；mean $\lambda$ 口径是错的 & 0.542342 /\allowbreak{} 0.406756 /\allowbreak{} 0.5 \checkmark \\
B & 二维显式$\cdot$r=0.2：实测 / 解析有限网格谱半径 & 0.9901506725 & 0.9901506724 \checkmark \\
B & 二维显式$\cdot$r=0.3：实测 / 解析有限网格谱半径 & 1.3852260087 & 1.3852260087 \checkmark \\
B & ADI$\cdot$r=0.3：实测 / 解析有限网格谱半径 & 0.9853345425 & 0.9853345424 \checkmark \\
B & ADI$\cdot$r=1.0：实测 / 解析有限网格谱半径 & 0.9519439461 & 0.9519439461 \checkmark \\
B & ADI$\cdot$r=10.0：实测 / 解析有限网格谱半径 & 0.8175781104 & 0.8175781101 \checkmark \\
B & ADI$\cdot$r=100.0：实测(幂迭代)/ 解析有限网格谱半径 & 0.9800770882 & 0.9799791594(\allowbreak{}低 9.99e-05) \checkmark \\
B & nx=20 的模态集合边界 & Dirichlet: (\allowbreak{}n$-$2)$\pi$/\allowbreak{}(\allowbreak{}n$-$1) < $\pi$；Neumann: $\pi$；周期: 2$\pi$$\cdot$floor(\allowbreak{}n/\allowbreak{}2)/\allowbreak{}n & 2.9762456718 /\allowbreak{} 3.1415926536 /\allowbreak{} 3.1415926536 \checkmark \\
B & nx=21 的模态集合边界 & Dirichlet: (\allowbreak{}n$-$2)$\pi$/\allowbreak{}(\allowbreak{}n$-$1) < $\pi$；Neumann: $\pi$；周期: 2$\pi$$\cdot$floor(\allowbreak{}n/\allowbreak{}2)/\allowbreak{}n & 2.9845130209 /\allowbreak{} 3.1415926536 /\allowbreak{} 2.9919930034 \checkmark \\
B & nx=41 的模态集合边界 & Dirichlet: (\allowbreak{}n$-$2)$\pi$/\allowbreak{}(\allowbreak{}n$-$1) < $\pi$；Neumann: $\pi$；周期: 2$\pi$$\cdot$floor(\allowbreak{}n/\allowbreak{}2)/\allowbreak{}n & 3.0630528373 /\allowbreak{} 3.1415926536 /\allowbreak{} 3.0649684425 \checkmark \\
\end{longtable}


\section{它会怎样骗你（本章版）}

\begin{enumerate}
  \item \textbf{拿"解看起来好不好"当稳定性判据}：$r=0.51$ 已越界，但 $\sin(\pi x)$
        初值下峰值 $0.6099$、误差 $3.20\times10^{-4}$（稳定的 $r=0.45$ 是 $2.64\times10^{-4}$），
        光滑得无可挑剔；同一个 $r$ 换成噪声初值，峰值 $42.88$、误差 $43.48$。
        更极端的 $r=0.6$：连正弦初值都被**舍入噪声**顶到 $759$（$|G|_{\max}=1.3963$，
        $134$ 步把 $10^{-16}$ 放大 $18$ 个数量级）。
  \item \textbf{把 $r\le\frac12$ 当成普适常数}：有限网格上 Dirichlet 是
        $0.5034331$（$n=20$）、$0.5030970$（$n=21$）；Neumann 恰是 $0.5000000002$；
        奇数周期网格是 $0.5028080$。三条曲线都收敛到 $\frac12$，没有一条等于 $\frac12$。
        二维是 $\frac14$（实测 $0.2515484898$）。
  \item \textbf{变系数时用平均系数代进判据}：$a(x)=1+x$ 上实测阈值 $s_*=0.5423415824$，
        平均系数口径给 $0.4068$（偏小 $25\%$），冻结系数判据给 $0.5$（保守 $8\%$）。
        "用 $\max$ 不用 $\mathrm{mean}$"是对的方向，"阈值恒等于 $\frac12$"不是。
  \item \textbf{把"无条件稳定"读成"随便取 $\Delta t$ 都准"}：BTCS 在 $r=50$ 的谱半径是
        $0.6461$（稳），但 $L_\infty$ 误差是 $r=0.4$ 时的 $86.2$ 倍；
        CN 在 $r=500$ 的谱半径 $0.9980$——稳定到"每步翻符号、永不衰减"。
        稳定只保证不发散。
  \item \textbf{拿连续解当靶子测时间阶}：CN 会被量成 $-0.0004$、BTCS 成 $0.4615$，
        而换半离散靶子是 $2.0000$ 与 $0.9994$。二维 ADI 同样：$-0.0018$ 对 $2.0000$。
        $\Delta t$ 缩小 $32$ 倍误差只从 $5.446\times10^{-4}$ 动到 $5.467\times10^{-4}$——
        量到的根本不是时间误差。
  \item \textbf{固定 $r$ 加密网格，然后宣称"一阶格式也是二阶"}：
        FTCS/BTCS/CN 在 $nx=11/21/41/81$ 上分别是 $2.0159/1.9977/2.0008$，
        因为 $r$ 固定时 $\Delta t=O(\Delta x^2)$，时间误差被压成同阶。
        这条路径测不出时间阶，换固定 $\Delta t$ 加密就现出一阶。
  \item \textbf{把幂迭代测到的数当谱半径的等号}：Neumann/周期下不减均值永远是
        $1.000000000000$（常数模态 $|G|=1$）；即使投影了，$nx=101$、$r=0.05$ 时顶部
        特征值挤在一起，实测只能到 $1.8\times10^{-4}$ 的相对精度，
        ADI 在 $r=100$ 时低估 $1.0\times10^{-4}$（$\lambda_2/\lambda_1=1-1.6\times10^{-4}$，
        $4000$ 步分不开）。容差要按谱间隙给。
  \item \textbf{把注释里的断言当成事实}：我原先写"Dirichlet 的恒等行会让幂迭代收敛到 $1$，
        所以也必须投影"——实测否定了它（投影前后都是 $0.996978987$，差 $5\times10^{-11}$）。
        原因是那两个等于 $1$ 的特征值对应的是**左**特征向量，而步进器每步都把端点钉成 $0$，
        迭代走不到那些方向。判据的方向感来自推导，数值来自测量。
  \item \textbf{模态分解时不管分母}：初值是纯 $\sin(\pi x)$ 时 $k\ge2$ 的系数精确为 $0$，
        直接相除得到 $22.1$、$18.5$ 这种"放大率"。分母是舍入级的时候，
        比值不是数据；本章给相对门槛（$10^{-8}\max|a_k|$）并置掩码。
  \item \textbf{只看分母，忘了分子也会被淹没}：$r=0.8$ 的 FTCS 在 $t=0.05$ 时长到
        $2.2^{100}\approx10^{34}$，其余模态的投影上叠着 $10^{18}$ 量级的舍入，
        "实测存活率"于是给出 $1.5\times10^{2}$ 的偏差。噪声尺度要取 $\|u\|_\infty$，
        不能取"投影出来这几个系数"的量级——按后者定标时 $8$ 个模态里有 $7$ 个被误判为可测。
        这一条是写帧的时候才发现的：$\texttt{meta.n\_absent}$ 与
        $\texttt{meta.n\_masked\_by\_roundoff}$ 现在分别记录两种原因。
  \item \textbf{以为余弦族对普通内积正交}：Neumann 网格上 $\cos(\pi x)$ 在
        $\cos(3\pi x)$ 上漏出 $0.0476=1/21$（$O(1/n)$ 的系统性偏差，不是舍入）。
        要用端点半权重（DCT-II 的正交关系）才对，修正后同一系数是 $1.2\times10^{-16}$。
  \item \textbf{周期网格上少补一个 $\Delta x$}：周期网格的 $x$ 到 $L-\Delta x$ 为止，
        $\cos(2\pi kx/L)$ 的周期要补一个 $\Delta x$ 才是 $L$。少补这一项，
        波数差 $2.5\%$，$L_\infty$ 误差从 $5.3\times10^{-4}$ 抬到 $3.56\times10^{-2}$——
        比格式自身的截断误差大两个数量级，看上去像"解算错了"。
  \item \textbf{把"残差小"当"解准"}：显式格式的回代残差**恰好是 $0.0$**（它根本没有解方程
        这一步），BTCS/CN 是 $6.2\times10^{-17}$、$7.9\times10^{-17}$；
        而同一批 BTCS 的解在 $r=50$ 时误差是 $r=0.4$ 的 $86.2$ 倍。
        残差量的是"是不是这个离散方程的解"，不是"这个离散方程对不对"。
  \item \textbf{把"总热量不变"当"解不变"}：Neumann/周期的漂移是 $3.2\times10^{-14}$ 与
        $3.4\times10^{-14}$（结构性的，与 $r$ 无关；显式格式只有 $2.2\times10^{-16}$），
        但常数模态永不衰减，数值解会冻结在一个与真解不同的常数上，误差不随时间减小；
        Dirichlet 的流失量 $0.247739$ 与解析值 $0.247975$ 一致到 $0.09\%$
        ——一个是舍入，一个是物理。
\end{enumerate}
